\section{Verification of \FenceDD}
\label{sec:mechanization}

The numerical theorems in \Cref{sec:numerics} concern the functional models
defined there, whereas users call the public C wrappers in
\Cref{sec:overview}. Following
earlier layered verifications of numerical C programs
\citep{AppelBertot2020,Kellison2023LAProof,Tekriwal2023Jacobi}, we prove three
obligations: kernel refinement, soundness of the runtime checks, and their
composition into each wrapper's public contract.  We illustrate these proofs
using division of a double word by a binary64 value; the other operations
follow the same structure.

\subsection{C Kernels Implement Their Functional Models}
\label{sec:kernel-refinement}

The numerical theorem in \Cref{thm:dw-div-fp-numerical} concerns the
functional model \DWDivFP{}, while the wrappers call the C kernel in
\Cref{fig:dw-div-fp-kernel}.  Suppose that \texttt{x} initially stores
$(x_h,x_\ell)$ and that \texttt{dw\_div\_fp\_kernel} returns with
\texttt{z} storing $(z_h,z_\ell)$.  Using VST, we prove
\begin{equation}
  (z_h,z_\ell)=\DWDivFP(x_h,x_\ell,y),
  \label{eq:dw-div-fp-kernel-refinement}
\end{equation}
where \DWDivFP{} is instantiated with binary64. The Rocq definition of
\DWDivFP{} uses the IEEE~754 types and operations provided by VCFloat2
\citep{AppelKellison2024}. In the VST proof, \texttt{FPCompCert} bridge
lemmas relate these operations to CompCert, while the specifications of the
transformation kernels account for the internal C calls.  This refinement identifies
the computed pair but requires no numerical condition and proves no error
bound.  We prove an analogous result for every arithmetic kernel in the
library.

\subsection{Runtime Checks Establish the Mathematical Conditions}
\label{sec:check-soundness}

To apply the numerical theorem to either wrapper for dividing a double-double
by a binary64 value, we must show that successful runtime checks establish its
mathematical hypotheses. The
following theorem connects the weak and strong checks in
\Cref{fig:dw-div-fp-input-check,fig:dw-div-fp-check} and the finiteness check
in \Cref{fig:finite-checks} to the conditions in
\Cref{sec:division-conditions}.

\begin{theorem}[Soundness of the Checks for Division by a Binary64 Value]
\label{thm:division-checks}
Suppose that \texttt{x} stores the binary64 pair $(x_h,x_\ell)$.

\begin{enumerate}
  \item If the weak check returns $1$, then
    $\mathcal{I}_{\mathrm{div}}(x_h,x_\ell,y)$ holds.

  \item If the strong check returns $1$, then
    $\mathcal{S}_{\mathrm{div}}(x_h,x_\ell,y)$ holds.
\end{enumerate}

If \texttt{z} stores $(z_h,z_\ell)$ and the finiteness check returns $1$,
then both components are finite.
\end{theorem}

The proofs account for short-circuit control flow using the specifications of
the finiteness functions, \texttt{fabs}, and \texttt{ldexp}.  The main
arithmetic argument derives the required bound on the exact quotient from the
rounded quotient comparison, using exact threshold representability and
monotonicity of rounding.  These results prove soundness, not completeness: a
failed runtime check need not imply that the corresponding mathematical
condition is false.

\subsection{Verified Contracts for Dividing a Double-Double by a Binary64
Value}
\label{sec:division-wrapper-contracts}

The wrapper proofs combine kernel refinement, check soundness, and the
numerical theorem.  \Cref{fig:dw-div-fp-vst-spec} shows the clauses of the VST
specification for the default wrapper in
\Cref{fig:default-division-wrapper} that state its numerical contract; routine
memory assertions are omitted.

\begin{figure}[t]
  \begin{lstlisting}[
    language=Coq,
    basicstyle=\ttfamily\normalsize,
    numbers=left
  ]
Definition dw_div_fp_spec :=
  DECLARE _dw_div_fp
  WITH z_ptr: val, x_ptr: val,
       x_pair: ftype Tdouble *
               ftype Tdouble,
       y: ftype Tdouble
  PRE [tptr dword, tptr dword, tdouble]
    PROP()
    ...
  POST [tint]
    EX status : int,
    PROP(default_wrapper_status
           status;
      status = Int.zero ->
      DWDivFP_guarantee
        (fst x_pair)
        (snd x_pair) y)
    RETURN(Vint status)
    SEP((if Int.eq status Int.zero
      then data_at Tsh dword
        (f2val
          (DWDivFP
            (fst x_pair)
            (snd x_pair) y))
        z_ptr
      else data_at_ Tsh dword z_ptr);
      ...).
  \end{lstlisting}
\caption{Relevant clauses of the VST specification for the default division
wrapper. On success, the postcondition supplies the numerical guarantee.
Omitted assertions describe the parameters and unchanged input object.}
  \label{fig:dw-div-fp-vst-spec}
\end{figure}


The empty \texttt{PROP} clause means that the caller supplies no numerical
hypothesis.  On normal return, the postcondition states that the output object
stores the binary64 result of \DWDivFP{} and that
\texttt{DWDivFP\_guarantee} holds.  In the Rocq
development, \texttt{DWDivFP\_guarantee} names the conclusion of
\Cref{thm:dw-div-fp-numerical}: the represented result satisfies the stated
error bound, and the returned pair satisfies the double-word relation.

For the default wrapper, the weak check establishes
$\mathcal{I}_{\mathrm{div}}$, kernel refinement identifies the result with
\DWDivFP{}, and the result check establishes $\mathcal{F}_{\mathrm{div}}$.
For the alternative wrapper, the strong check establishes
$\mathcal{S}_{\mathrm{div}}$ and hence result finiteness.  These are the two
alternatives in \Cref{cor:dw-div-conditions}, so both wrappers provide the same
contract.

\subsection{Compilation and the Verification Boundary}
\label{sec:verification-boundary}

The C proofs assume that \texttt{double} is binary64 and that arithmetic uses
round-to-nearest, ties-to-even, with gradual underflow.
Compilation and execution must preserve these semantics without reassociating
floating-point expressions or flushing subnormal values to zero.
Floating-point exceptions are nontrapping, which allows the default wrapper to
reach the result check on Line~9 of \Cref{fig:default-division-wrapper} when
the kernel produces an infinity or NaN \citep{IEEE754-2019}.
The verified functions call \texttt{fma}, \texttt{fabs}, \texttt{ldexp}, and
\texttt{exit}.  In the proofs, we use the behavior described by their VST
specifications, but the platform implementations of these functions and the
floating-point hardware are outside the verification boundary
\citep{Appel2023VSTlib}.

We prove the C-level contracts using VST.  For a CompCert build, we apply
semantic preservation to carry these results to assembly
\citep{Leroy2009,Boldo2013CompCert}.  Compilation with another
compiler is not covered by this result and must independently preserve the
required IEEE~754 semantics.
